\chapter{Làm nhóm, ăn theo và đạo đức rất linh hoạt}
\markboth{Làm nhóm, ăn theo và đạo đức rất linh hoạt}{Làm nhóm, ăn theo và đạo đức rất linh hoạt}

\section{Em không seen tin nhắn, em đang detox nhóm lớp}
\begin{truyen}{Em không seen tin nhắn, em đang detox nhóm lớp}{Classroom Talk}
\chuhoa{M}{ột bạn bị hỏi lý do mất hút liền đáp: ``Dạo này em detox điện thoại, nên nhóm nhắn em không đọc được.''}

Thầy hỏi: ``Detox kỹ vậy mà ảnh story với tin crush em vẫn thả tim đúng giờ, còn deadline nhóm thì chết lặng?''

Bạn ấy cười quê: ``Cái đó... do em lỡ thấy.''

Thầy nói: ``Ừ. Nhiều người không hề mất kết nối. Họ chỉ chọn mất kết nối đúng phần có trách nhiệm thôi.''

\textit{(Selective silence is rarely detox. It is usually avoidance wearing a cleaner label.)}
\end{truyen}

\begin{baihoc}
Mất sóng đúng lúc phải làm việc là kỹ năng rất giỏi... nếu mục tiêu là làm người khác phát điên.
\textit{(Going mysteriously offline at work time is a powerful skill... if your goal is to drive others crazy.)}
\end{baihoc}

\begin{ghinhoanh}
\textbf{Short English to remember:}

Selective silence is still avoidance.

Responsibility does not disappear because notifications are muted.
\end{ghinhoanh}

\begin{tuhocnhanh}
1. Em đang quá tải thật, hay đang né việc nhóm? \textit{(Are you truly overwhelmed, or avoiding group work?)}

2. Nếu em là người chờ phản hồi, em sẽ chịu nổi kiểu biến mất đó không? \textit{(If you were waiting for a reply, would you tolerate that disappearance?)}
\end{tuhocnhanh}
\ngancach

\section{Em không làm slide vì em mạnh chiến lược}
\begin{truyen}{Em không làm slide vì em mạnh chiến lược}{Classroom Talk}
\chuhoa{C}{ó bạn rất tự tin: ``Em không giỏi làm nội dung tay chân. Em thiên về chiến lược nên em đứng ngoài nhìn tổng thể cho nhóm.''}

Thầy hỏi: ``Tổng thể đó có tạo được một cái file, một trang slide, hay ít nhất một đoạn nội dung nào chưa?''

Bạn ấy nói nhỏ: ``Dạ em định góp ý ở cuối.''

Thầy chốt: ``Chiến lược mà không kéo theo đóng góp thật thì thường chỉ là du lịch tư duy trên công sức người khác.''

\textit{(Claiming the strategic role means little when there is no concrete work behind it.)}
\end{truyen}

\begin{baihoc}
Không phải cứ đứng ngoài chỉ tay là thành người có tầm nhìn. Nhiều khi chỉ là chưa chịu làm phần nặng tay.
\textit{(Standing aside and giving directions does not automatically make you visionary. Sometimes it only shows reluctance to do real work.)}
\end{baihoc}

\begin{ghinhoanh}
\textbf{Short English to remember:}

Strategy without contribution is cheap.

Real vision still carries work.
\end{ghinhoanh}

\begin{tuhocnhanh}
1. Em có đang tự tô vai trò của mình cho oai hơn thực tế không? \textit{(Are you inflating your role to sound more important than it is?)}

2. Nhóm sẽ mất gì nếu ai cũng “làm chiến lược” như em? \textit{(What happens if everyone in the team does strategy the way you do?)}
\end{tuhocnhanh}
\ngancach

\section{Em đến trễ nhưng em đem năng lượng}
\begin{truyen}{Em đến trễ nhưng em đem năng lượng}{Classroom Talk}
\chuhoa{M}{ột bạn vào muộn gần hết buổi họp nhóm rồi cười tươi: ``Em tới trễ nhưng em đem năng lượng tích cực cho mọi người.''}

Thầy hỏi: ``Năng lượng đó có gánh được phần việc ba bạn kia vừa làm không?''

Bạn ấy đáp: ``Ít ra em cũng làm không khí vui mà thầy.''

Thầy nói: ``Nhóm đang cần người kéo việc, không cần thêm mascot phát sáng lúc việc gần xong.''

\textit{(Cheerfulness is welcome, but it cannot replace labor when everyone else has already carried the load.)}
\end{truyen}

\begin{baihoc}
Tinh thần tốt là phụ kiện đẹp. Phần việc làm ra mới là thân bài của trách nhiệm.
\textit{(Good energy is a nice accessory. Actual work is the body of responsibility.)}
\end{baihoc}

\begin{ghinhoanh}
\textbf{Short English to remember:}

Positive energy is nice.

It is not a substitute for effort.
\end{ghinhoanh}

\begin{tuhocnhanh}
1. Em đang góp giá trị thật hay chỉ góp bầu không khí? \textit{(Are you contributing real value, or just mood?)}

2. Nếu ai cũng đến muộn nhưng rất vui vẻ, nhóm còn cái gì để nộp? \textit{(If everyone arrives late but cheerful, what would the group still have to submit?)}
\end{tuhocnhanh}
\ngancach

\section{Em không ăn theo, em chỉ biết hòa hợp}
\begin{truyen}{Em không ăn theo, em chỉ biết hòa hợp}{Classroom Talk}
\chuhoa{C}{ó bạn ký tên vào bài nhóm rất tự nhiên rồi nói: ``Em đâu có ăn theo. Em chỉ hòa hợp với thành quả chung của tập thể thôi.''}

Thầy hỏi: ``Hòa hợp kiểu gì mà công người khác chạy mồ hôi, còn tên em đáp xuống đúng lúc file hoàn tất?''

Bạn ấy chống chế: ``Thì nhóm là của chung mà thầy.''

Thầy đáp: ``Ừ. Của chung không có nghĩa là ai cũng được quyền làm ký sinh trùng học đường một cách thơ mộng.''

\textit{(Teamwork means shared ownership, not shared permission to contribute nothing and still collect credit.)}
\end{truyen}

\begin{baihoc}
Tinh thần tập thể không phải cái chăn để đắp lên sự ăn theo cho đỡ lạnh lương tâm.
\textit{(Team spirit is not a blanket to cover freeloading and keep the conscience warm.)}
\end{baihoc}

\begin{ghinhoanh}
\textbf{Short English to remember:}

Shared credit needs shared effort.

Harmony is not freeloading with better vocabulary.
\end{ghinhoanh}

\begin{tuhocnhanh}
1. Em có đang hưởng phần điểm mình chưa góp đủ công không? \textit{(Are you receiving credit you did not really earn?)}

2. Nếu nhóm chấm công công bằng tuyệt đối, em được mấy phần? \textit{(If the team graded contribution fairly, what share would you deserve?)}
\end{tuhocnhanh}
\ngancach